% --- Άσκηση 33895 (33895.tex) ---
\Askhsh{33895}{4}
\begin{ylh}{2}
    \item 6.1. Η Έννοια της Συνάρτησης
    \item 6.2. Γραφική Παράσταση Συνάρτησης
    \item 6.3. Η Συνάρτηση $f(x)=ax+\beta$
\end{ylh}
Δίνεται η συνάρτηση $f(x) = \dfrac{4x^2 - 2(a + 3)x + 3a}{2x - 3}$, με παράμετρο $a \in \mathbb{R}$.

\begin{alist}
\item Να βρείτε το πεδίο ορισμού της $f$. \monades{5}
\item Να αποδείξετε ότι $f(x) = 2x - a$, για κάθε $x$ που ανήκει στο πεδίο ορισμού της $f$. \monades{8}
\item Να βρείτε την τιμή του $a \in \mathbb{R}$, αν η γραφική παράσταση της $f$ διέρχεται από το σημείο $(1, -1)$. \monades{7}
\item Να βρείτε, αν υπάρχουν, τα σημεία τομής της γραφικής παράστασης της $f$ με τους άξονες $x'x$ και $y'y$. \monades{5}
\end{alist}
